%% 03_methods.tex
\section{Methods}
\label{sec:methods}

\subsection{Dataset}
\label{sec:methods:dataset}

\paragraph{Base population.}
We generate 200 synthetic patients using Synthea~\citep{walonoski2018synthea}
with a fixed random seed, then apply a custom specialty-regimen overlay that assigns each
patient to one of three complexity tiers:
\begin{itemize}
  \item Tier 1 (Low): single long-acting injectable on a fixed recurring interval
    (e.g.\ every 8 weeks, modelling LAIs such as paliperidone palmitate~\citep{correll2024pp6m}).
  \item Tier 2 (Medium): cyclic regimen (on-for-$N$-weeks / off-for-$M$-weeks)
    requiring cycle-position reasoning.
  \item Tier 3 (High): multi-drug regimen (biologic primary + oral adjunct + PRN
    rescue) with staggered or interdependent schedules, motivating cross-resource reasoning.
\end{itemize}
The final distribution is 63 Tier-1, 63 Tier-2, and 74 Tier-3 patients.  All 200 FHIR bundles
are validated against a reference R4B pre-write conformance validator
(\texttt{mh\_integration/r4b\_validator.py}) using \texttt{fhir.resources} $\geq$8.0; all
200 bundles pass with no mandatory element failures.

\paragraph{Paired representations.}
Each patient is represented in two parallel forms derived from the same FHIR bundle:
\begin{enumerate}
  \item FHIR bundle (structured): Synthea base plus the specialty overlay
    (MedicationRequest, MedicationAdministration, Medication, CarePlan, Encounter, Patient,
    Practitioner), frozen at git tag \texttt{dataset-freeze-v1} (commit \texttt{98959f8};
    manifest SHA-256 \texttt{4d0da93e\ldots}).
  \item LLM narrative: generated from the bundle by Llama~3.3-70B via Groq with
    a fixed system prompt (hash \texttt{518bc71d87b7}), fixed temperature (0.0), and a
    mid-window administration date instruction to surface temporal entities.  A deterministic
    templated-narrative ablation (no LLM) provides an information-ceiling baseline.
\end{enumerate}

\paragraph{Fidelity audit.}
Narrative quality is verified by an independent programmatic audit that recovers eight entity
classes (medication name, regimen start/end, tier integer, tier description, administration
dates, administration count, dosage schedule) directly from FHIR resources (not from the
generation prompt, fixing the circular-audit problem).  After one prompt-iteration step to
add mid-window administration dates, both generators reached 100\% weighted entity
recall across all 200 patients and all 2,403 checks (O1 gate: PASS).  The same audit
was replicated on the templated-narrative ablation with the same result.

\paragraph{Question bank.}
We generate 13,800 questions (200 patients $\times$ 69 templates per patient).  Each
question record additionally carries 2 paraphrase variants of the question prompt;
both paraphrases share a single ground-truth answer and are scored as one row.  The
69 templates span five PSP adherence-indicator families:
\begin{enumerate}
  \item Next-dose lookup (\texttt{temporal\_lookup}, $n=3{,}200$): ``What is the
    next scheduled administration on reference date $d$?''
  \item Dose-history aggregation (\texttt{regimen\_aggregation}, $n=2{,}400$): ``How many
    doses have been administered in the past 90 days?'' (MPR numerator)
  \item Coverage-window reasoning (\texttt{regimen\_compliance}, $n=4{,}200$): ``What
    proportion of the past 90 days was the patient covered?'' (PDC)
  \item Missed-dose detection (\texttt{temporal\_comparison}, $n=2{,}400$): ``Is there
    evidence of a missed or delayed dose?''
  \item Persistence (\texttt{cross\_resource}, $n=1{,}600$): ``Is the patient still on
    therapy?'' (cross-resource component questions)
\end{enumerate}
All ground-truth values are computed programmatically from FHIR bundles using pure-Python
adherence metric functions (MPR, PDC, is\_persistent, consecutive\_missed\_doses); no manual
annotation and no LLM-as-judge is used anywhere in evaluation.  The question bank SHA-256
is \texttt{2216f58e\ldots} (post-W2 MPR-window and trivial-template corrections;
see the Reproducibility section).

\subsection{Retrieval Systems}
\label{sec:methods:systems}

All three systems share controlled variables (Table~\ref{tab:controlled}):
embedding model \texttt{BAAI/bge-large-en-v1.5} (local via sentence-transformers),
answer LLM \texttt{qwen/qwen3-32b} via Groq (temperature 0.0, single run per question),
ChromaDB local vector store, $k=5$ retrieved chunks, chunk size 500 tokens, overlap 50 tokens.

\begin{table}[!htbp]
\centering
\caption{Controlled variables held constant across Systems A, B, C.}
\label{tab:controlled}
\begin{tabular}{ll}
\toprule
Parameter & Value \\
\midrule
Embedding model & \texttt{BAAI/bge-large-en-v1.5} (local) \\
Answer LLM & \texttt{qwen/qwen3-32b} via Groq Developer plan \\
Temperature & 0.0 \\
Runs per question & 1 \\
Top-$k$ & 5 \\
Chunk tokens & 500 \\
Chunk overlap & 50 tokens \\
Vector store & ChromaDB (local) \\
Answer max tokens & 512 tokens \\
\bottomrule
\end{tabular}
\end{table}

\paragraph{System A (Narrative RAG).}
LLM narratives are split into 500-token chunks with 50-token overlap, embedded, and indexed.
At query time the top-5 chunks are retrieved by cosine similarity and concatenated into the
answer prompt.

\paragraph{System B (Structured RAG, Naive).}
Each FHIR resource in the bundle is serialised to a flat text representation, chunked,
embedded, and indexed.  Retrieval is identical to System~A.  Mean tokens-in per call: 4,463
vs.\ System~A's 490, a 9$\times$ context expansion from serialising full FHIR JSON.

\paragraph{System C (Structured RAG, Resource-Aware).}
Three retrieval primitives are applied before dense search:
(1) \emph{typed resource filter}: the question family is mapped to a resource-type allowlist
(e.g.\ MedicationAdministration + CarePlan for dose-history questions);
(2) \emph{reference-chain traversal}: MedicationRequest references are expanded to
Medication, Practitioner, and Patient resources;
(3) \emph{temporal pre-filter}: only resources with dates intersecting a window around
the question's reference date are retained.  Dense search then selects top-5 from the
filtered subset.  Mean tokens-in: 2,175 (51\% less than System~B).

\subsection{Evaluation}
\label{sec:methods:eval}

\paragraph{Accuracy.}
The scorer (\texttt{eval/score.py}) applies deterministic type-detection rules to classify
each ground truth as date, numeric, boolean/categorical, list, or free-text, then applies
type-specific scoring (exact-match binary; 2\% relative tolerance for floats; Jaccard
coefficient for lists).  The primary metric is exact-match rate.  Partial credit is reported as a secondary metric.

\paragraph{Retrieval recall@k.}
For Systems~B and~C only (System~A's narrative chunks carry no FHIR resource IDs), recall@$k$
is computed by checking whether the top-$k$ retrieved chunk IDs include the patient's primary
MedicationRequest or specialty CarePlan (a conservative v1 heuristic; see Limitations).

\paragraph{Statistical tests.}
Paired bootstrap (10,000 resamples, paired by question\_id, seed 42) produces 95\%
confidence intervals for pairwise accuracy differences.  McNemar's test (uncorrected
$\chi^2$, 1~df) is applied to exact-match for each pair.  Wilson score 95\%~CIs are reported
for all per-cell proportions.  Three planned comparisons (A vs.\ B, A vs.\ C, B vs.\ C);
no multiple-comparison correction (pre-specified confirmatory design).

\paragraph{Error taxonomy.}
For each system, 50 failures are drawn stratified by question family (10 per family, seed 42)
and classified by deterministic regex rules into five categories (four primary failure modes plus a residual ``Other''):
(1) temporal-anchor failure,
(2) reasoning-truncated,
(3) N/A misuse (phantom value or spurious N/A),
(4) wrong-entity / enum error,
(5) other.
Single-rater (TJ); no LLM-based triangulation was performed. The taxonomy is fully deterministic and reproducible from the regex rules in \texttt{analysis/run\_error\_taxonomy.py} (decision C3).

\paragraph{Feature-extraction arm (O6).}
Three patient-level feature sets are constructed.
FS-Structured consists of 19 FHIR-derived numeric features: MPR-90d/180d/full, PDC-90d/full,
gap statistics, event counts, days-since-last-dose, tier one-hot, resource-type counts, and age.
FS-Narrative consists of 6 heuristic regex features on LLM narratives: date-mention count,
miss-word count, dose-count mentions, tier-label mentions, narrative length, and negation-near-dose
count (no LLM scoring).
FS-Aware consists of 8 per-patient aggregates from System~C retrieval traces.
Two classifiers (LightGBM with $n\_\text{estimators}=200$, $\text{learning\_rate}=0.05$;
logistic regression with $C=1.0$, class\_weight=balanced) are evaluated via stratified 5-fold
cross-validation.  Paired bootstrap AUC CIs are computed at the patient level from
out-of-fold predictions.

\paragraph{API and cost disclosure.}
Narrative generation (200 calls) used the Groq free tier.  The full evaluation matrix
(approximately 41,400 calls) used the Groq Developer paid plan (authorised 2026-05-05;
hard ceiling \$20~USD).  Cost computed from token totals at list-price rates over live
(non-cached) calls only: System~A \$4.74, System~B \$21.07, System~C \$6.88 (total
\$32.68).  System~B exceeds the \$20 single-system ceiling; the overrun was accepted
per the decision log entry of 2026-05-06.  System~C remains within budget.  Actual
billed amounts were lower than the list-price total because of response-cache reuse
on reproducible re-runs.
